\clearpage\thispagestyle{empty}
\pdfbookmark[0]{Quellen zu Pallas (25. April 1812)}{pallas1812apr25-sources}
\begin{center}\textbf{Quellen zu Pallas (25. April 1812)}\end{center}
\noindent Diese KI-Übertragung ins heutige Deutsch folgt den \href{https://gdz.sub.uni-goettingen.de/id/PPN236007467}{\textit{Werken}, Band~VI} (Göttingen, 1874), S.~349--351. Der Bericht erschien unter der Überschrift Göttingen in den \textit{Göttingischen gelehrten Anzeigen}, Heft~67 vom 25. April 1812, S.~\href{https://api.digitale-sammlungen.de/iiif/image/v2/bsb10538681_00701/full/1600,/0/default.jpg}{657}, \href{https://api.digitale-sammlungen.de/iiif/image/v2/bsb10538681_00702/full/1600,/0/default.jpg}{658}, \href{https://api.digitale-sammlungen.de/iiif/image/v2/bsb10538681_00703/full/1600,/0/default.jpg}{659} und \href{https://api.digitale-sammlungen.de/iiif/image/v2/bsb10538681_00704/full/1600,/0/default.jpg}{660}. Das digitalisierte Exemplar stammt aus der Bayerischen Staatsbibliothek. Der Bericht enthält zwei Mitteilungen aus einem Brief Friedrich Wilhelm Bessels an Gauß vom 26. März.
\par\smallskip\noindent In der Zeitschrift steht auf S.~659 bei der Deklination von Sequens zuerst die Zahl 37, und für die jährliche Bewegung des Hauptsterns in Deklination stehen dort 3 Sekunden 321. In der Werkausgabe stehen 47 und 2 Sekunden 321, mit positivem Vorzeichen bei der jährlichen Bewegung. Die Übertragung folgt der Werkausgabe. KI-Übertragung und KI-Satz.
\clearpage
